\chapter*{Notation}
\addcontentsline{toc}{chapter}{Notation}

\def\arraystretch{1.5}
\section*{Numbers and Arrays}
\begin{longtable}{p{0.25\textwidth}p{0.65\textwidth}}
$\displaystyle a$ & A scalar (integer or real)\\
$\displaystyle \va$ & A vector\\
$\displaystyle \vsa{1:N}$ & A sequence of vectors from $1$ to $N$, such that 
$\displaystyle \vsa{1:N} = (\va_1,\dots,\va_N)$
\\
$\displaystyle \mA$ & A matrix\\
$\displaystyle \tA$ & A tensor\\
$\displaystyle \mI_n$ & Identity matrix with $n$ rows and $n$ columns\\
$\displaystyle \mI$ & Identity matrix with dimensionality implied by context\\
$\displaystyle \ve^{(i)}$ & Standard basis vector $[0,\dots,0,1,0,\dots,0]$ with a 1 at position $i$\\
$\displaystyle \text{diag}(\va)$ & A square, diagonal matrix with diagonal entries given by $\va$\\
$\displaystyle \ra$ & A scalar random variable\\
$\displaystyle \rva$ & A vector-valued random variable\\
$\displaystyle \rmA$ & A matrix-valued random variable\\
\end{longtable}

%
\section*{Sets and Graphs}
\def\arraystretch{1.5}
\begin{longtable}{p{0.25\textwidth}p{0.65\textwidth}}
$\displaystyle \sA$ & A set\\
$\displaystyle \R$ & The set of real numbers \\
$\displaystyle \mathbb{C}$ & The set of complex numbers \\
% NOTE: do not use \R^+, because it is ambiguous whether:
% - It includes 0
% - It includes only real numbers, or also infinity.
% We usually do not include infinity, so we may explicitly write
% [0, \infty) to include 0
% (0, \infty) to not include 0
$\displaystyle \{0, 1\}$ & The set containing 0 and 1 \\
$\displaystyle \{0, 1, \dots, n \}$ & The set of all integers between $0$ and $n$\\
$\displaystyle [a, b]$ & The real interval including $a$ and $b$\\
$\displaystyle (a, b]$ & The real interval excluding $a$ but including $b$\\
$\displaystyle \sA \backslash \sB$ & Set subtraction, i.e., the set containing the elements of $\sA$ that are not in $\sB$\\
\end{longtable}
%
%
\section*{Indexing}
\def\arraystretch{1.5}
\begin{longtable}{p{0.25\textwidth}p{0.65\textwidth}}
$\displaystyle \eva_i$ & Element $i$ of vector $\va$, with indexing starting at 1.\\
$\displaystyle \eva_{-i}$ & All elements of vector $\va$ except for element $i$ \\
$\displaystyle \emA_{i,j}$ & Element $i, j$ of matrix $\mA$ \\
$\displaystyle \mA_{i, :}$ & Row $i$ of matrix $\mA$ \\
$\displaystyle \mA_{:, i}$ & Column $i$ of matrix $\mA$ \\
$\displaystyle \etA_{i, j, k}$ & Element $(i, j, k)$ of a 3-D tensor $\tA$\\
$\displaystyle \tA_{:, :, i}$ & 2-D slice of a 3-D tensor\\
$\displaystyle \erva_i$ & Element $i$ of the random vector $\rva$ \\
\end{longtable}
%
%
\section*{Linear Algebra Operations}
\def\arraystretch{1.5}
\begin{longtable}{p{0.25\textwidth}p{0.65\textwidth}}
$\displaystyle \mA^\top$ & Transpose of matrix $\mA$ \\
% Wikipedia uses \circ for element-wise multiplication but this could be confused with function composition
\end{longtable}
%
\section*{Calculus}
\def\arraystretch{1.5}
\begin{longtable}{p{0.25\textwidth}p{0.65\textwidth}}
$\displaystyle\frac{d y} {d x}$ & Derivative of $y$ with respect to $x$\\ [2ex]
$\displaystyle \frac{\partial y} {\partial x} $ & Partial derivative of $y$ with respect to $x$ \\
$\displaystyle \nabla_\vx y $ & Gradient of $y$ with respect to $\vx$ \\
$\displaystyle \frac{\partial f}{\partial \vx} $ & Jacobian matrix $\mJ \in \R^{m\times n}$ of $f: \R^n \rightarrow \R^m$\\
$\displaystyle \int f(\vx) d\vx $ & Definite integral over the entire domain of $\vx$ \\
\end{longtable}
%
%
%
%
\section*{Probability and Information Theory}
\def\arraystretch{1.5}
\begin{longtable}{p{0.25\textwidth}p{0.65\textwidth}}
$\displaystyle \ra \bot \rb$ & The random variables $\ra$ and $\rb$ are independent\\
$\displaystyle P(\ra)$ & A probability distribution over a discrete variable\\
$\displaystyle p(\ra)$ & A probability distribution over a continuous variable, or over
a variable whose type has not been specified\\
$\displaystyle \ra \sim P$ & Random variable $\ra$ has distribution $P$\\% so thing on left of \sim should always be a random variable, with name beginning with \r
$\displaystyle  \E_{\rx\sim P} [ f(x) ]$ & Expectation of $f(x)$ with respect to $P(\rx)$ \\
$\displaystyle  Q_{\tau, \rx\sim P} [ f(x) ]$ & Quantile of $f(x)$ with respect to $P(\rx)$ with coverage $\tau$ \\
$\displaystyle \mathcal{N} ( \vx ; \vmu , \mSigma)$ & Gaussian distribution %
over $\vx$ with mean $\vmu$ and covariance $\mSigma$ \\
\end{longtable}
%
\section*{Functions}
\def\arraystretch{1.5}
\begin{longtable}{p{0.25\textwidth}p{0.65\textwidth}}
$\displaystyle f: \sA \rightarrow \sB$ & The function $f$ with domain $\sA$ and range $\sB$\\
$\displaystyle f \circ g $ & Composition of the functions $f$ and $g$ \\
  $\displaystyle f(\vx ; \vtheta) $ & A function of $\vx$ parametrized by $\vtheta$.
  (Sometimes we write $f(\vx)$ and omit the argument $\vtheta$ to lighten notation) \\
$\displaystyle \log x$ & Natural logarithm of $x$ \\
$\displaystyle || \vx ||_p $ & $\normlp$ norm of $\vx$ \\
$\displaystyle || \vx || $ & $\normltwo$ norm of $\vx$ \\
$\displaystyle x^+$ & Positive part of $x$, i.e., $\max(0,x)$\\
\end{longtable}

%
\section*{Datasets and Distributions}
\def\arraystretch{1.5}
\begin{longtable}{p{0.25\textwidth}p{0.65\textwidth}}
$\displaystyle \pdata$ & The data generating distribution\\
$\displaystyle \ptrain$ & The empirical distribution defined by the training set\\
$\displaystyle \sX$ & A set of training examples\\
$\displaystyle \vx^{(i)}$ & The $i$-th example (input) from a dataset\\
$\displaystyle y^{(i)}\text{ or }\vy^{(i)}$ & The target associated with $\vx^{(i)}$ for supervised learning\\
$\displaystyle \mX$ & The $m \times n$ matrix with input example $\vx^{(i)}$ in row $\mX_{i,:}$\\
\end{longtable}